\sect{मार्गशीर्ष-शुक्ल-मोक्षदा-एकादशी-माहात्म्यम्}
\label{sec:padma-margashirsha-shukla-mokshada}

\ganapatyadiDhyanam
\padmaPuranam

\dnsub{अथ एकचत्वारिंशोऽध्यायः॥४१}

\uvacha{युधिष्ठिर उवाच}

\twolineshloka
{वन्दे विष्णुं विभुं साक्षाल्लोकत्रयसुखावहम्}
{विश्वेशं विश्वकर्तारं पुराणं पुरुषोत्तमम्}% ॥१॥

\twolineshloka
{पृच्छामि देवदेवेश संशयोऽस्ति महान्मम}
{लोकानां च हितार्थाय पापानां क्षयहेतवे}% ॥२॥

\threelineshloka
{मार्गशीर्षे सिते पक्षे भवेदेकादशी तु या}
{किं नाम को विधिस्तस्याः को देवस्तत्र पूज्यते}
{एतदाचक्ष्व मे स्वामिन् विस्तरेण यथातथम्}% ॥३॥

\uvacha{श्रीकृष्ण उवाच}

\twolineshloka
{साधु पृष्टं त्वया राजन् साधु ते विमलं यशः}
{कथयिष्यामि राजेन्द्र हरिवासरमुत्तमम्}% ॥४॥

\twolineshloka
{उत्पन्ना चासिते पक्षे द्वादशी मम वल्लभा}
{मार्गशीर्षोत्पत्तिरिति मम देहसमुद्भवा}% ॥५॥

\twolineshloka
{सुरासुरवधार्थाय ह्युत्पन्ना भरतर्षभ}
{कथिता च मया सा वै तवाग्रे राजसत्तम}% ॥६॥

\twolineshloka
{पूर्वा चैकादशी राजंस्त्रैलोक्ये सचराचरे}
{मार्गशीर्षेऽसिते पक्षे उत्पत्तिरिति नामतः}% ॥७॥

\twolineshloka
{अतः परं प्रवक्ष्यामि मार्गशीर्षे सिता तु या}
{यस्याः श्रवणमात्रेण वाजपेयफलं लभेत्}% ॥८॥

\twolineshloka
{मोक्षा नामेति सा प्रोक्ता सर्वपापहरा परा}
{देवं दामोदरं राजन् पूजयेच्च प्रयत्नतः}% ॥९॥

\twolineshloka
{तुलस्या मञ्जरीभिश्च धूपैर्दीपैः प्रयत्नतः}
{पूर्वेण विधिना चैव दशम्येकादशी तथा}% ॥१०॥

\twolineshloka
{मोक्षा चैकादशी नाम्ना महापातकनाशिनी}
{रात्रौ जागरणं कार्यं नृत्यगीतस्तवैर्मम}% ॥११॥

\twolineshloka
{शृणु राजन् प्रवक्ष्यामि दिव्यां पौराणिकीकथाम्}
{यस्याः श्रवणमात्रेण सर्वपापक्षयो भवेत्}% ॥१२॥

\twolineshloka
{अधोयोनिगतश्चैव पितरो यस्य पापतः}
{अस्याश्च पुण्यदानेन मोक्षं यान्ति न संशयः}% ॥१३॥

\twolineshloka
{चम्पके नगरे रम्ये वैष्णवैश्च विभूषिते}
{वैखानसो नाम नृपः पुत्रवत् पालयेत् प्रजाः}% ॥१४॥

\twolineshloka
{वसन्ति बहवो विप्रा वेदवेदाङ्गपारगाः}
{ऋद्धिमत्यः प्रजास्तस्य राज्ञो वैखानसस्य हि}% ॥१५॥

\twolineshloka
{एवं राज्यं प्रकुर्वाणो रात्रौ स्वप्नस्य मध्यतः}
{स्वकीयपितरो दृष्ट्वा अधोयोनिगता नृप}% ॥१६॥

\twolineshloka
{एवं दृष्ट्वा च तान् सर्वान् विस्मयाविष्टमानसः}
{कथयामास वृत्तान्तं स्वप्नजातं द्विजान् प्रति}% ॥१७॥

\uvacha{राजोवाच}

\twolineshloka
{मया स्वपितरो दृष्ट्वा नरकोपगता द्विजाः}
{तारयेति तनूज त्वमस्मान्निरयसागरात्}% ॥१८॥

\twolineshloka
{एवं ब्रुवाणास्ते नूनं रोदमाना मुहुर्मुहुः}
{मया दृष्टा द्विजश्रेष्ठा एतस्माच्च न मे सुखम्}% ॥१९॥

\twolineshloka
{एतद् राज्यं मम महत् सुखदायि न विद्यते}
{अश्वा गजास्तथा सर्वे रोचन्ते मे न भो द्विजाः}% ॥२०॥

\twolineshloka
{न दारा न सुता मह्यं रोचन्ते द्विजसत्तमाः}
{किं करोमि क्व गच्छामि हृदयं मेऽवरुध्यते}% ॥२१॥

\twolineshloka
{तद् व्रतं तं तपोयोगं येनैव मम पूर्वजाः}
{मोक्षं प्रयान्ति सद्यो वै कथ्यतां च द्विजोत्तमाः}% ॥२२॥

\twolineshloka
{पुत्रे तु जीवितप्राये बलीयसि महात्मनि}
{पिताऽस्ति नरके घोरे तस्य पुत्रस्य किं फलम्}% ॥२३॥

\uvacha{ब्राह्मणा ऊचुः}

\twolineshloka
{पर्वतस्य मुने राजन् निकटे चाऽऽश्रमो महान्}
{गम्यतां राजशार्दूल भूतं भव्यं विजानतः}% ॥२४॥

\twolineshloka
{तेषां श्रुत्वा ततो वाक्यं राजा वैखानसो महान्}
{जगाम चाऽऽशु तत्रैव चाऽऽश्रमं पर्वतस्य च}% ॥२५॥

\twolineshloka
{ब्राह्मणैर्वेष्टितो राजा राजभिश्च समन्वितः}
{आश्रमं विपुलं तस्य सम्प्राप्तो राजसत्तमः}% ॥२६॥

\twolineshloka
{तत्रर्ग्वेद-यजुर्वेद-सामाध्ययन-कोविदैः}
{वेष्टितं मुनिभिश्चैव द्वितीयं ब्रह्मणो यथा}% ॥२७॥

\twolineshloka
{दृष्ट्वा तं मुनिशार्दूलं राजा वैखानसस्तथा}
{दण्डवत् प्रणतिं कृत्वा पस्पर्श चरणौ मुनेः}% ॥२८॥

\twolineshloka
{पप्रच्छ कुशलं तस्य सप्तस्वङ्गेष्वसौ मुनिः}
{राज्ये निष्कण्टकत्वं च राज्ञः सौख्यसमन्वितम्}% ॥२९॥

\uvacha{राजोवाच}

\twolineshloka
{तव प्रसादाद् भो स्वामिन् कुशलं मेऽङ्गसप्तसु}
{भक्ता ये विष्णुविप्रेषु कथं तेषां च विघ्नता}% ॥३०॥

\twolineshloka
{मया स्वपितरो दृष्टाः स्वप्ने च नरके स्थिताः}
{कस्य पुण्यस्य सामर्थ्यान्मोक्षं यान्ति द्विजोत्तम}% ॥३१॥

\twolineshloka
{अयं मे संशयः स्वामिन् प्रष्टुं तं त्वामुपागतः}
{उपायः कश्चिदेवात्र कर्तव्यो मुनिसत्तम}% ॥३२॥

\twolineshloka
{एतद् वाक्यं ततः श्रुत्वा पर्वतो मुनिसत्तमः}
{ध्यानस्तिमितनेत्रोऽभूत् तपस्वी ब्रह्मसन्निभः}% ॥३३॥

\twolineshloka
{मुहूर्तमेकं ध्यानस्थो भूपतिं प्रत्युवाच ह}
{ज्ञातं हि तव राजेन्द्र पितॄणां पूर्वचेष्टितम्}% ॥३४॥

\twolineshloka
{पूर्वजन्मनि तातस्ते क्षत्रियो राज्यगर्वितः}
{सपत्न्या ऋतुकाले तु राजधर्मप्रवर्तितः}% ॥३५॥

\twolineshloka
{गतो ग्रामे तु तां त्यक्त्वा कार्यार्थी निजयोषितम्}
{तव पित्रा तु तस्याश्च न दत्तमृतुदानकम्}% ॥३६॥

\twolineshloka
{तेन पापप्रभावेन नरके पितृभिः सह}
{पतितो राजशार्दूल तव तातः सुदारुणे}% ॥३७॥

\twolineshloka
{ततः पुनरुवाचेदं राजा वैखानसो मुनिम्}
{केन व्रतप्रभावेन मोक्षस्तेषां भवेन्मुने}% ॥३८॥

\uvacha{मुनिरुवाच}

\twolineshloka
{मार्गशीर्षे सिते पक्षे मोक्षा नामेति नामतः}
{सर्वैश्चेतद् व्रतं कार्यं पित्रे पुण्यं प्रदीयताम्}% ॥३९॥

\twolineshloka
{तेन पुण्यप्रभावेन मोक्षस्तेषां भविष्यति}
{सत्यमेतन्महाभाग ब्रह्मणो वचनं यथा}% ॥४०॥

\twolineshloka
{मुनेर्वाक्यं ततः श्रुत्वा स्वगृहं पुनरागतः}
{मार्गशीर्षस्तथा मासः प्राप्तः कष्टेन तेन वै}% ॥४१॥

\twolineshloka
{मुनेर्वाक्येन तत्कृत्वा व्रतं वैखानसो नृपः}
{अददत् पुण्यमखिलैः सार्धं पित्रे स भूमिपः}% ॥४२॥

\twolineshloka
{दत्ते पुण्यक्षणेनैव पुष्पवृष्टिरभूद्दिवि}
{वैखानसस्य तातो वै पितृभिर्मोक्षमाविशत्}% ॥४३॥

\twolineshloka
{राजानं चान्तरिक्षे स गिरं पुण्यामुवाच ह}
{स्वस्ति स्वस्तीति ते पुत्र प्रोच्य चैवं दिवं गतः}% ॥४४॥

\twolineshloka
{एवं यः कुरुते राजन् मोक्षामेकादशीं शुभाम्}
{तस्य पापानि नश्यन्ति मृतो मोक्षमवाप्नुयात्}% ॥४५॥

\twolineshloka
{नातः परतरा काचिन्मोक्षदैकादशी भवेत्}
{पुण्यसङ्ख्यां न जानामि राजन् मे प्रियकृद् व्रतम्}% ॥४६॥

\twolineshloka
{चिन्तामणिसमा ह्येषा नृणां मोक्षप्रदायिनी}
{पठनाच्छ्रवणादस्या वाजपेयफलं लभेत्}% ॥४७॥

॥इति श्रीपाद्मे महापुराणे पञ्चपञ्चाशत्साहस्र्यां संहितायामुत्तरखण्डे उमापतिनारदसंवादान्तर्गतकृष्णयुधिष्ठिरसंवादे मार्गशीर्ष-शुक्ल-मोक्षदा-एकादशी-माहात्म्यं नाम एकचत्वारिंशोऽध्यायः॥४१॥

\genMangalaShloka

\hyperref[sec:ekadashi_mahatmyam_padma_puranam]{\closesub}
\clearpage

